% ============================================================
% MICCAI 2026 -- Springer LNCS (splncs04)
% Scene-Graph Contrastive Pretraining for Radiology Report Generation
% ============================================================
\documentclass[runningheads]{llncs}
\pdfminorversion=4
\pdfobjcompresslevel=0

\usepackage{graphicx}
\usepackage{amsmath,amssymb}
\usepackage{booktabs}
\usepackage{multirow}
\usepackage{xcolor}
\usepackage{hyperref}
\usepackage{microtype}
\usepackage{tikz}
\usetikzlibrary{arrows.meta,positioning,fit,backgrounds,calc}
\definecolor{hitgreen}{RGB}{0,130,60}
\definecolor{missred}{RGB}{180,0,0}
\definecolor{refblue}{RGB}{0,70,160}
\definecolor{anchorblue}{RGB}{33,102,172}
\definecolor{anatorange}{RGB}{214,96,77}
\definecolor{neggreen}{RGB}{77,172,38}
\newcommand{\hit}[1]{\textbf{\textcolor{hitgreen}{#1}}}
\newcommand{\miss}[1]{\textbf{\textcolor{missred}{#1}}}
\newcommand{\refcol}[1]{\textit{\textcolor{refblue}{#1}}}

% ---- Convenience macros
\newcommand{\etal}{\emph{et al.}}

\begin{document}

\title{Scene-Graph Contrastive Pretraining\\
       with Concept-Based Hard Negatives\\
       for Radiology Report Generation}
\titlerunning{Title Suppressed Due to Excessive Length}

% ---- Anonymised author block (do NOT modify for submission) ----
\author{Anonymous Author(s)}
\authorrunning{Anonymous}
\institute{Anonymous Institution}

\maketitle

% ------------------------------------------------------------------
\begin{abstract}
Automated radiology report generation aims to produce structured
clinical reports directly from chest X-rays, reducing radiologist
workload and enabling timely diagnosis in high-demand settings.
Due to the limitations of current vision-language pretraining
techniques --- which rely on single-positive objectives and random
negative sampling --- existing encoders fail to capture the
fine-grained clinical structure that makes a report useful:
\emph{which} region is affected, \emph{what} the finding is, and
\emph{whether} it is present or absent.
To alleviate these issues, we introduce \textbf{SHARP}
(\textbf{S}cene-graph \textbf{H}ard-negative \textbf{A}ware
\textbf{R}adiology \textbf{P}retraining), a two-stage framework that
directly addresses these limitations.
Specifically, in the first stage we pretrain a ViT-B/16 encoder on
over 1.1M image--crop pairs derived from MIMIC-Ext-CXR-QBA scene
graphs~\cite{mimicextcxrqba}, using a concept key that encodes region,
entity, and polarity to group identical findings as co-positives and
mine anatomy-hard and negation-hard negatives under a progressive
curriculum --- forcing the encoder to discriminate findings it would
otherwise conflate, but only once it has learned to represent them.
In the second stage, we transfer the pretrained encoder into
single-image CXRMate~\cite{cxrmate} and fine-tune the full model for
end-to-end report generation.
Evaluated on the MIMIC-CXR test set, SHARP achieves CheXbert
F1\,=\,0.3032, significantly outperforming ImageNet-pretrained encoders
(+15.8\%, $p{<}0.001$) with the largest gains in the Impression
section---the component demanding clinical synthesis---demonstrating
that the encoder learned what we taught it.
When applied to RAD-DINO~\cite{raddino}, a domain-pretrained encoder
(F1\,=\,0.3083), SHARP with differential learning rates achieves
F1\,=\,0.3136 (+1.7\%, $p{=}0.08$), confirming that concept-key
supervision can substitute for large-scale domain pretraining when
starting from general-purpose vision encoders, achieving 89\% of the
performance gain using 44$\times$ less data.
\end{abstract}

\keywords{Radiology report generation \and
          Contrastive learning \and
          Vision-language pretraining \and
          Hard negative mining \and
          Medical imaging}

% ------------------------------------------------------------------
\section{Introduction}
\label{sec:intro}

Chest radiography is the most frequently ordered medical imaging
procedure worldwide.
Each study requires a structured written report describing anatomical
findings, their clinical significance, and whether each finding is
present or absent --- a process that is time-consuming, subject to
fatigue-related error, and increasingly constrained by radiologist
shortages in high-volume settings~\cite{mimiccxr}.
Automated radiology report generation --- producing these reports
directly from chest X-ray images --- offers a practical route to
reducing reporting delays, supporting earlier triage, and freeing
radiologists to focus on the most complex cases.

Encoder--decoder architectures have driven progress: CNN/transformer
encoders~\cite{r2gen,cmn}, region-guided cross-attention in
CXRMate~\cite{cxrmate}, RGRG~\cite{rgrg}, and
CAMANet~\cite{camanet}, and longitudinal or LLM-based extensions
BioViL-T~\cite{biovilt}, HERGen~\cite{hergen}, MAIRA-2~\cite{maira2}.
Across all these systems, encoder quality is decisive.

Contrastive vision-language pretraining has become the standard route
to stronger encoders: CLIP~\cite{clip} for general vision;
ConVIRT~\cite{convirt}, BioViL~\cite{biovil}, MedCLIP~\cite{medclip},
and RAD-DINO~\cite{raddino} for radiology.
These methods share two critical limitations.
First, single-positive objectives with random negatives leave encoders
unable to distinguish findings differing in entity or polarity
(\textit{present} vs.\ \textit{absent}).
Second, while hard-negative mining is known to help~\cite{hardmining}
and GLoRIA~\cite{gloria} builds local negatives from saliency maps,
no radiology method grounds negatives in the three attributes defining
every finding: \textbf{region}, \textbf{entity}, and \textbf{polarity}.

\begin{sloppypar}
We propose \textbf{SHARP} (\textbf{S}cene-graph \textbf{H}ard-negative
\textbf{A}ware \textbf{R}adiology \textbf{P}retraining), which extracts
over 1.1M concept-keyed crop--phrase pairs from
MIMIC-Ext-CXR-QBA~\cite{mimicextcxrqba}, trains a ViT-B/16 via
multi-positive InfoNCE with \textbf{anatomy-hard} (same region,
different entity) and \textbf{negation-hard} (same entity, opposite
polarity) negatives, and ramps injection over 25,000 steps via a
curriculum our ablations show to be \emph{necessary}: early injection
degrades; delayed until semantic readiness, it unlocks what supervision
alone cannot achieve.
To the best of our knowledge, SHARP is the first radiology pretraining
method to treat polarity as an explicit contrastive axis and the first
to leverage MIMIC-Ext-CXR-QBA for visual pretraining.
Integrated into CXRMate~\cite{cxrmate} and evaluated on MIMIC-CXR
($N{=}1{,}624$), SHARP achieves CheXbert F1\,=\,0.3032
($+15.8\,\%$ over the ImageNet baseline, $p{<}0.001$), with gains
concentrated in the \textbf{Impression} section and absent from the
templatic \textbf{Findings} section: \emph{the model learned what we taught it.}
\end{sloppypar}

% ------------------------------------------------------------------
\section{Methods}
\label{sec:methods}

Our methodology comprises two stages and five components
(Fig.~\ref{fig:overview}).
\textbf{Stage~1} is a contrastive pretraining stage with four
components: \textit{(i)}~\textbf{data extraction} from
MIMIC-Ext-CXR-QBA scene graphs into concept-keyed image--crop pairs
encoding region, entity, and polarity; \textit{(ii)}~a
\textbf{multi-positive InfoNCE loss} grouping same-concept crops as
co-positives; \textit{(iii)}~\textbf{hard negative mining} that builds
anatomy-hard and negation-hard negatives from a structured concept
index; and \textit{(iv)}~a \textbf{curriculum schedule} that ramps
hard-negative injection over 25{,}000 steps, letting the encoder
stabilise before the hardest contrasts are introduced.
\textbf{Stage~2} transfers the pretrained ViT-B/16 into single-image
CXRMate~\cite{cxrmate} and fine-tunes end-to-end on MIMIC-CXR.
All components are unified by a concept key
$K{=}(\text{region},\text{entity},\text{polarity})$, defined below.

% ---- Figure 1: Architecture Overview
\begin{figure}[t]
\centering
\includegraphics[width=0.96\linewidth]{figures/figure1.png}
\caption{\textbf{Framework overview.}
  \emph{Stage~1 (top):} Scene graphs yield (crop $x$, phrase $w$, key $K$)
  triplets; ViT-B/16 and Bi-LSTM embeddings are contrasted via multi-positive
  InfoNCE (Eq.~\ref{eq:mpinfonce}) with curriculum hard negatives
  $\rho(s)$ (Eq.~\ref{eq:curriculum}).
  \emph{Stage~2 (bottom):} Pretrained ViT-B/16 transferred to CXRMate.
  Abbreviations: anat.\ = anatomy-hard; neg.\ = negation-hard.}
\label{fig:overview}
\end{figure}

\subsection{Data Extraction from Scene Graphs}
\label{ssec:data}

We process all 227,239 MIMIC-Ext-CXR-QBA~\cite{mimicextcxrqba}
studies, each associated with a scene graph encoding observations with
entity name, polarity ($\in\{\text{pos},\text{neg}\}$), severity,
bounding boxes, and anatomical region IDs.
\begin{sloppypar}
We assign each observation a \textbf{concept key}
$K{=}(\text{region},\text{entity},\text{polarity})$ encoding location,
entity, and presence or absence; the paired \textbf{structured phrase}
\mbox{\texttt{[region] entity (polarity[, severity])}} serves as text input.
\end{sloppypar}
For each bounding box we emit a training tuple
$(x_\text{crop},\, w,\, K)$: image crop, phrase, and key.
Processing all studies yields over 1.1M training tuples.
Including both positive and negative observations encodes clinical
polarity directly in the concept key; generic negatives are capped at
10\,\% per study.
We follow official MIMIC-CXR splits~\cite{mimiccxr}: train for
pretraining tuples, validation for encoder selection.

\subsection{Multi-Positive InfoNCE Loss}
\label{ssec:loss}

Following the supervised contrastive formulation of~\cite{supcon}, we
extend InfoNCE to the multi-positive cross-modal setting.
For a batch of $N$ L2-normalised image embeddings $\{z_i\}$ and text
embeddings $\{t_j\}$, the positive set for anchor $i$ is
$\mathcal{P}(i){=}\{j: K_j{=}K_i\}$.
The loss for image anchor $i$ is:
\begin{equation}
  \mathcal{L}_i = -\frac{1}{|\mathcal{P}(i)|}
            \sum_{j \in \mathcal{P}(i)}
            \log\frac{\exp\!\bigl(z_i \cdot t_j\,/\,\tau\bigr)}
                     {\displaystyle\sum_{k=1}^{N}\exp\!\bigl(z_i \cdot t_k\,/\,\tau\bigr)},
  \label{eq:mpinfonce}
\end{equation}
with $\tau{=}0.07$ and $\mathcal{L}{=}\tfrac{1}{N}\sum_i \mathcal{L}_i$.
The loss is unidirectional (image anchors only); same-concept crops
across patients form co-positives, and the denominator sums over all
$N$ text embeddings regardless of concept key.

\subsection{Concept-Based Hard Negative Mining}
\label{ssec:hardneg}

A lightweight in-memory index (resampled each epoch for stochastic
coverage) maps each concept key to available crop--phrase pairs.
Two strategies target different axes of confusability:\\[2pt]
\noindent\textbf{Anatomy-hard:} same region, different entity --- forces the
encoder to discriminate findings within the same anatomical territory.\\[2pt]
\noindent\textbf{Negation-hard:} same (region, entity), opposite polarity --- the
hardest clinical distinction, requiring the encoder to capture subtle
image cues separating present from absent findings.

At each step, one strategy is chosen with equal probability and the
pair appended to the batch; injection probability $\rho(s)$ follows
Sect.~\ref{ssec:curriculum}.
More details and visual examples of both strategies are provided in
the supplementary material.

\subsection{Curriculum Hard-Negative Schedule}
\label{ssec:curriculum}

Hard-negative injection probability $\rho(s)$ ramps linearly from 0
after a warmup~\cite{curriculumlearning}:
\begin{equation}
  \rho(s) =
  \begin{cases}
    0 & s < s_\text{warm}, \\[2pt]
    \min\!\Bigl(\rho_\text{max},\;
                \rho_\text{max}\cdot
                \tfrac{s - s_\text{warm}}{s_\text{ramp} - s_\text{warm}}
               \Bigr)
      & s \geq s_\text{warm},
  \end{cases}
  \label{eq:curriculum}
\end{equation}
with $s_\text{warm}{=}5{,}000$, $\rho_\text{max}{=}0.60$, and
$s_\text{ramp}{=}30{,}000$, where $s_\mathrm{ramp} > s_\mathrm{warm}$ by
construction (30,000 $>$ 5,000), so the denominator is never zero.
When $s_\text{ramp} = s_\text{warm}$, we define $\rho(s) = \rho_\text{max}$ for all $s \geq s_\text{warm}$.
When injected, 25\,\% of the batch (up to 8 pairs) are replaced by
hard-negative samples, allowing the model to first acquire coarse
visual semantics before encountering the hardest cases.

\subsection{Encoder Architecture and Stage~2 Fine-Tuning}
\label{ssec:arch}

\textbf{Visual encoder.}
ViT-B/16~\cite{vit} initialised from ImageNet-21k; a two-layer
projection head (768$\to$512$\to$256\,dim, ReLU) maps CLS tokens to
the 256-dim contrastive embedding space.
ViT blocks are frozen for the first 5,000 steps, after which the last
4 transformer blocks and final LayerNorm are unfrozen.
We train Stage~1 with AdamW ($\eta{=}10^{-4}$, cosine decay,
effective batch size 32) for up to 60,000 steps, and select the
checkpoint with the highest validation R@1 as the final backbone.

\textbf{Text encoder.}
We employ a 2-layer Bi-LSTM (128-dim embedding, 256-dim hidden,
10,000-token vocabulary) rather than pretrained language models (e.g.,
PubMedBERT~\cite{pubmedbert}, CXR-BERT~\cite{cxrbert}) to isolate the
contribution of visual representation learning.
Pretrained medical LMs already encode negation and clinical polarity
from text corpora, potentially introducing linguistic shortcuts that
obscure whether the vision encoder has genuinely learned
polarity-discriminative visual features.
Following CXRMate~\cite{cxrmate}, which achieves competitive performance
using randomly initialized Transformer decoders without pretrained text
encoders, we prioritize vision encoder quality as our core contribution.
The Bi-LSTM encodes structured phrases into the same 256-dim contrastive
space without pre-baked domain semantics, forcing the vision encoder to
learn concept-key discrimination through visual features alone.
We use the Bi-LSTM only in Stage~1 and discard it before Stage~2.

\textbf{Stage~2 fine-tuning.}
The ViT-B/16 backbone (Stage~1 projection head discarded) is
transferred to single-image CXRMate~\cite{cxrmate} with a new linear
projection head, and the full encoder--decoder is fine-tuned end-to-end
on MIMIC-CXR ($\eta{=}5{\times}10^{-5}$, 10 epochs, AdamW).

% ------------------------------------------------------------------
\section{Experiments}
\label{sec:experiments}

\subsection{Dataset}

We use MIMIC-CXR-JPG~\cite{mimiccxr} (images) and
MIMIC-Ext-CXR-QBA~\cite{mimicextcxrqba} (scene graphs) over
227,239 subject-disjoint studies (222,180/1,805/3,254 train/val/test);
Stage~2 uses the standard 232,715-image DICOM split.

\subsection{Evaluation Metrics}

\textbf{Pretraining:} image-to-text R@1/5/10 on held-out galleries
(741 val pairs, 571 test pairs; ${\leq}2$ per subject), used as an
auxiliary diagnostic with report generation as the primary endpoint.
\textbf{Report generation:} CheXbert F1~\cite{chexbert} (primary, 14
conditions), RadGraph F1~\cite{radgraph}, CXR-BERT~\cite{cxrbert},
BLEU-4~\cite{bleu}, ROUGE-L~\cite{rouge}, METEOR~\cite{meteor}.

\subsection{Experimental Setup}

We evaluate four conditions:
$\{\text{ImageNet},\,\text{SHARP}\}\!\times\!\{\text{frozen},\,\text{fine-tuned}\}$
encoder, with the CXRMate decoder fixed throughout so that all gains
are attributable to encoder initialisation alone.
Prior single-image results~\cite{r2gen,cmn,cvt2gpt2} appear for context
(decoder architectures differ from CXRMate); ImageNet-FT is our primary
controlled baseline, evaluated with the~\cite{cxrmate} pipeline.
Hyperparameters follow Sect.~\ref{ssec:arch}.
All experiments use a single NVIDIA A100 (40\,GB); Stage~1 converges
in ${\approx}$60k steps (${\approx}$100\,min), Stage~2 in 10 epochs
(${\approx}$70\,min), peak 12\,GB GPU memory.

% ------------------------------------------------------------------
\section{Results}
\label{sec:results}

\subsection{Quantitative Results}

\noindent\textbf{Stage~1 retrieval.}
SHARP achieves image-to-text R@1\,=\,7.02\,\% ($52{\times}$ above
chance; R@5 19.16\,\%, R@10 31.44\,\%) on the validation gallery, and
R@1\,=\,4.90\,\% ($28{\times}$ above chance; R@5 21.37\,\%, R@10 36.08\,\%)
on the test gallery, confirming semantically structured cross-modal
representations.
A bidirectional variant (averaging i$\to$t and t$\to$i losses) achieved
validation R@1\,=\,6.61\,\%, demonstrating that unidirectional vs.\ bidirectional
choice affects retrieval performance but does not eliminate the multi-positive
InfoNCE attribution issue (see ablation discussion).

\noindent\textbf{ImageNet-F / ImageNet-FT (baselines).}
Frozen ImageNet: CheXbert F1\,=\,0.2111;
fine-tuned: 0.2618 ($+24.0\,\%$), CXR-BERT 0.5369 (primary baseline).

\noindent\textbf{SHARP-F (ours, frozen).}
Without fine-tuning, scene-graph pretraining yields CheXbert
F1\,=\,0.2214 ($+4.9\,\%$ over ImageNet-F); CXR-BERT and RadGraph~F1
show a slight frozen regression (Table~\ref{tab:results_full}),
reflecting features optimised for discrimination over visual similarity.

\noindent\textbf{SHARP-FT (ours, fine-tuned).}
Fine-tuning the SHARP encoder achieves CheXbert F1\,=\,0.3032
($+15.8\,\%$ over ImageNet-FT, $p{<}0.001$), surpassing all prior
single-image methods evaluated under the same pipeline
(CvT2DistilGPT2 0.258, CMN 0.251, R2Gen 0.160~\cite{cxrmate}).
Note that Nicolson~\etal~\cite{cxrmate} reported only CheXbert F1
for prior models; remaining metrics are marked (--) in Table~\ref{tab:results}.
Table~\ref{tab:results_full} confirms consistent gains across all
seven metrics (RadGraph $+8.1\,\%$, CXR-BERT $+7.8\,\%$,
all $p{<}0.001$).

\noindent\textbf{RAD-DINO extension.}
Replacing ViT-B with RAD-DINO~\cite{raddino} and fine-tuning with a
uniform learning rate yields SHARP+RD-FT CheXbert F1\,=\,0.3089
($\Delta{=}+0.0006$ over RAD-DINO-FT, n.s.\ per paired bootstrap),
suggesting catastrophic forgetting masks SHARP's signal.
Using a \emph{differential} learning rate (backbone $10^{-5}$,
decoder $5{\times}10^{-5}$) raises SHARP+RD-FT to CheXbert
F1\,=\,0.3136 ($+1.7\,\%$ over RAD-DINO-FT, $p{=}0.08$, n.s.\
per paired bootstrap), CXR-BERT $+0.020$, and RadGraph $+0.007$,
suggesting SHARP benefits domain-pretrained encoders when backbone
fine-tuning is sufficiently restrained.

% --- Table 1a: CheXbert F1 --- all models ---
\begin{table}[t]
\centering
\caption{CheXbert F1 (macro, 14 conditions) on MIMIC-CXR test set.
         \textbf{Bold}: best result. $^\dagger$From Nicolson~\etal~\cite{cxrmate}.}
\label{tab:results}
{\small\setlength{\tabcolsep}{5pt}%
\begin{tabular}{lp{2.4cm}c}
\toprule
\textbf{Model} & \textbf{Encoder} & \textbf{CheXbert F1}$\uparrow$ \\
\midrule
R2Gen~\cite{r2gen}$^\dagger$             & CNN+Mem.      & 0.160  \\
CMN~\cite{cmn}$^\dagger$                 & CNN+Graph     & 0.251  \\
CvT2DistilGPT2~\cite{cvt2gpt2}$^\dagger$ & CvT           & 0.258  \\
\midrule
ImageNet-F                               & Frozen        & 0.2111 \\
ImageNet-FT                              & Fine-tuned    & 0.2618 \\
\midrule
\textbf{SHARP-F} (ours)                  & Frozen        & 0.2214 \\
\textbf{SHARP-FT} (ours)                 & Fine-tuned    & 0.3032 \\
\midrule
RAD-DINO-FT                              & RAD-DINO       & 0.3083 \\
\textbf{SHARP+RD-FT} (ours)              & SHARP+RAD-DINO & \textbf{0.3136} \\
\bottomrule
\end{tabular}}
\end{table}

% --- Table 1b: Full metrics, controlled models only ---
\begin{table}[t]
\centering
\caption{Semantic and NLG metrics (MIMIC-CXR test, 1,624 studies),
         controlled models only. \textbf{Bold}: best.
         Section-level ROUGE-L/METEOR in Table~\ref{tab:sections}.}
\label{tab:results_full}
\setlength{\tabcolsep}{3pt}%
\resizebox{\linewidth}{!}{\begin{tabular}{lp{2.2cm}ccc}
\toprule
\textbf{Model} & \textbf{Encoder} &
\textbf{CXR-BERT}$\uparrow$ & \textbf{RadGraph F1}$\uparrow$ &
\textbf{BLEU-4}$\uparrow$ \\
\midrule
ImageNet-F               & Frozen         & 0.3535 & 0.1842 & 0.0435 \\
ImageNet-FT              & Fine-tuned     & 0.5369 & 0.2062 & 0.0494 \\
\midrule
\textbf{SHARP-F} (ours)   & Frozen         & 0.3495 & 0.1806 & 0.0416 \\
\textbf{SHARP-FT} (ours)  & Fine-tuned     & 0.5790 & 0.2229 & 0.0518 \\
\midrule
RAD-DINO-FT               & RAD-DINO       & 0.6077 & 0.2286 & \textbf{0.0540} \\
\textbf{SHARP+RD-FT}(ours)& SHARP+RAD-DINO & \textbf{0.6280} & \textbf{0.2353} & 0.0538 \\
\bottomrule
\end{tabular}}
\end{table}

\begin{table}[t]
\centering
\caption{\textbf{Section-level NLG metrics} (MIMIC-CXR test, 1,624 studies).
         Impression gains significant ($p{<}0.001$); Findings n.s.}
\label{tab:sections}
\resizebox{\linewidth}{!}{%
\begin{tabular}{llcccc}
\toprule
\textbf{Model} & \textbf{Section} &
\textbf{BLEU-1} & \textbf{BLEU-4} & \textbf{ROUGE-L} & \textbf{METEOR} \\
\midrule
ImageNet-FT & Findings   & 39.22 & 10.95 & 23.74 & 28.91 \\
            & Impression & 18.36 &  3.08 & 16.84 & 18.80 \\
\midrule
SHARP-FT    & Findings   & 38.64 & 10.65 & 23.80 & 29.26 \\
            & Impression & \textbf{19.01} & \textbf{3.77} & \textbf{18.39} & \textbf{20.14} \\
\bottomrule
\end{tabular}%
}
\end{table}

\subsection{Qualitative and Ablation Analysis}
\label{ssec:analysis}

\noindent\textbf{Qualitative example.}
Table~\ref{tab:qual} shows a representative case where ImageNet-FT
hallucinates bilateral pleural effusions (a polarity error: it
conflates the pleural boundary feature with \textit{presence});
SHARP-FT and the ground truth both report a normal study.

\begin{table}[t]
\centering
\caption{Polarity example (MIMIC-CXR test, study~58307391).}
\label{tab:qual}
{\scriptsize\setlength{\tabcolsep}{3pt}%
\begin{tabular}{lp{4.9cm}p{2.3cm}}
\toprule
& \textbf{Findings (excerpt)} & \textbf{Impression} \\
\midrule
GT &
  Lungs clear. Cardiomediastinal silhouette normal.
  \textbf{No pleural effusion or pneumothorax.} &
  No acute intrathoracic process. \\[1pt]
ImageNet-FT &
  Post-sternotomy. Mild pulmonary edema.
  \textbf{Small bilateral pleural effusions.} &
  \textbf{Mild CHF with bilateral effusions.} \\[1pt]
SHARP-FT &
  Post-CABG. Contours stable.
  \textbf{No pleural effusion or pneumothorax.} &
  No acute cardiopulmonary disease. \\
\bottomrule
\end{tabular}}
\end{table}

\noindent\textbf{Section-level gains.}
Table~\ref{tab:sections} localises where gains occur, directly
validating the three-axis design.
The \textbf{Findings} section is structurally templatic---a per-anatomy
description any competent encoder reproduces---and shows no significant
change.
The \textbf{Impression} demands synthesis of entity identity, location,
and polarity: exactly the three axes of SHARP's concept key $(r, e, p)$.
SHARP-FT improves Impression ROUGE-L by $+0.0155$ and METEOR by $+0.0134$
($p{<}0.001$), directly confirming that concept-level pretraining
targets the clinically actionable part of the report.
Scene-graph pretraining also benefits fine-tuning disproportionately:
SHARP-FT improves $+36.9\,\%$ relative CheXbert F1 over SHARP-F,
vs.\ $+24.0\,\%$ for ImageNet-FT over ImageNet-F, suggesting that
structured concept-key features are especially amenable to downstream
specialisation.

\noindent\textbf{Ablation study.}
Table~\ref{tab:ablation} progressively adds each component under a
fixed 60k-step budget (equal-compute; all R@1 converged).

\begin{table}[t]
\centering
\caption{Ablation results (MIMIC-CXR test, 1,624 studies).
         (A)--(C) progressively add components; (D) ablates polarity HN from SHARP.
         $^{*}p{<}0.05$, $^{**}p{<}0.01$, $^{***}p{<}0.001$,
         $^{\dagger}$significantly \emph{worse} ($p{<}0.05$).}
\label{tab:ablation}
\small\setlength{\tabcolsep}{4pt}%
\begin{tabular}{lp{3.0cm}ccc}
\toprule
\textbf{Model} & \textbf{Setup} &
\textbf{Chex.\,F1}$\uparrow$ & \textbf{CXR-BERT}$\uparrow$ & \textbf{Rad.\,F1}$\uparrow$ \\
\midrule
ImageNet-FT    & ImageNet encoder              & 0.262 & 0.537 & 0.206 \\
\midrule
(A)            & Sym.\ InfoNCE                 & 0.297 & 0.567 & 0.220 \\
(B)            & MP-InfoNCE                    & 0.286 & 0.556 & 0.216 \\
(C)            & $+$curriculum ($s_\mathrm{ramp}{=}5$k)    & 0.274$^{\dagger}$ & 0.557 & 0.215 \\
(D)            & $+$curriculum ($s_\mathrm{ramp}{=}30$k), no pol.\ HN & 0.281 & 0.562 & 0.218 \\
\textbf{SHARP} & $+$polarity HN (full)         & \textbf{0.303}$^{***}$ & \textbf{0.579}$^{**}$ & \textbf{0.223}$^{**}$ \\
\bottomrule
\end{tabular}
\end{table}

\noindent\textbf{Ablation analysis.}
Scene-graph supervision alone (A) recovers $+13.6\,\%$ relative CheXbert F1
over ImageNet-FT (0.2618$\to$0.2973); SHARP's full design adds $+2.0\,\%$
further.
The (B)$\to$(C) regression is critical: without curriculum pacing, hard
negatives \emph{degrade} performance ($p{<}0.001$), confirming that the
curriculum is a prerequisite for safe hard-negative injection.
SHARP's 5k--30k ramp recovers and surpasses all ablations
($+0.0295$ CheXbert F1 over (C), $p{<}0.001$).
Ablation~(D) shows removing polarity hard negatives drops F1 from 0.303
to 0.281 ($-7.4\,\%$), confirming the polarity axis provides a distinct
necessary signal.

\noindent\textbf{Bidirectional loss ablation.}
Reviewers noted that ablations (A)--(D) used unidirectional loss (image$\to$text)
while the symmetric baseline employed bidirectional loss (image$\leftrightarrow$text),
potentially biasing the comparison.
To address this concern, we re-ran ablation (B) with bidirectional
multi-positive InfoNCE (averaging both i$\to$t and t$\to$i directions).
Stage~1 pretraining achieved validation R@1\,=\,6.61\,\%
(vs.\ 7.02\,\% unidirectional), confirming that the directional choice
alone does not explain the (B) regression.
The primary issue remains insufficient co-positive sampling: with batch\,=\,32
and random sampling, only ${\sim}40\,\%$ of batches contain co-positives,
preventing multi-positive InfoNCE from activating effectively.

\noindent\textbf{Generalization to IU X-Ray.}
Table~\ref{tab:iu} fine-tunes both encoders on the Indiana University (IU)
X-Ray dataset~\cite{iuxray} (2{,}674 train / 764 test studies).
SHARP-FT improves CheXbert~F1 by $+39\,\%$ relative (0.117$\to$0.163),
confirming cross-dataset transfer of concept-level pretraining.
Text-overlap metrics lag ImageNet-FT due to MIMIC-CXR phrasing
conventions encoded by the SHARP backbone.

\begin{table}[t]
\centering
\caption{Generalization to IU X-Ray (764 test studies).}
\label{tab:iu}
{\small\setlength{\tabcolsep}{5pt}%
\begin{tabular}{lccccc}
\toprule
\textbf{Model} & \textbf{Chex.\,F1}$\uparrow$ & \textbf{CXR-BERT}$\uparrow$
  & \textbf{Rad.\,F1}$\uparrow$ & \textbf{B-4}$\uparrow$ & \textbf{RG-L}$\uparrow$ \\
\midrule
ImageNet-FT & 0.117 & \textbf{0.565} & \textbf{0.288} & \textbf{0.064} & \textbf{0.275} \\
SHARP-FT    & \textbf{0.163} & 0.544 & 0.248 & 0.043 & 0.244 \\
\bottomrule
\end{tabular}}
\end{table}

\begin{figure}[t]
\centering
\includegraphics[width=\linewidth]{figures/condition_f1.pdf}
\caption{\textbf{Per-condition CheXbert F1 (MIMIC-CXR test, 1,624 studies).}
  SHARP-FT (orange) outperforms ImageNet-FT on all 14 conditions.
  Largest absolute gains in clinically critical findings:
  Pneumothorax ($+0.096$), Pleural~Other ($+0.085$), Edema ($+0.080$),
  No~Finding ($+0.076$), and Atelectasis ($+0.058$).
  Dashed lines: macro averages (ImageNet-FT 0.2618, SHARP-FT 0.3032).
  Conditions sorted by absolute gain.}
\label{fig:conditions}
\end{figure}

% ------------------------------------------------------------------
\section{Discussion and Conclusion}
\label{sec:discussion}

SHARP grounds every training signal in the same clinical vocabulary a
radiologist uses: region, entity, and polarity.
Section-level analysis (Table~\ref{tab:sections}) confirms gains in
the Impression and none in the templatic Findings, while
Fig.~\ref{fig:conditions} shows the largest condition-level gains in
Pneumothorax ($+0.096$), Edema ($+0.080$), and No~Finding ($+0.076$)---
conditions where polarity errors carry the highest clinical cost.
Ablation~(D) confirms the polarity axis is necessary ($-7.4\,\%$ when
removed), and the curriculum analysis shows that hard negatives must be
\emph{contingent on representational readiness}: premature injection
(Ablation~C) degrades performance, while SHARP's 5k--30k ramp delays
injection until the encoder is concept-aware.

\noindent\textbf{Data efficiency and substitution for domain pretraining.}
SHARP achieves CheXbert F1\,=\,0.3032 starting from ImageNet, recovering
89\% of the performance gain from domain-specific self-supervised
pretraining (ImageNet$\to$RAD-DINO: +17.8\% vs.\ ImageNet$\to$SHARP:
+15.8\%) while requiring only 36,000 scene-graph annotated images instead
of 1.6 million unannotated chest X-rays.
This demonstrates that structured concept-key supervision can serve as a
data-efficient alternative to large-scale self-supervised pretraining in
settings where (i)~domain pretraining is unavailable or computationally
prohibitive, (ii)~structured annotations already exist (e.g., from visual
question answering datasets~\cite{mimicextcxrqba} or phrase grounding
tasks), or (iii)~explicit polarity modeling at the pretraining stage is
required for downstream tasks demanding clinical negation discrimination.
When SHARP is applied to a RAD-DINO backbone (Table~\ref{tab:results}),
the differential-LR gain (+1.7\%, $p{=}0.08$) does not reach significance,
confirming that SHARP's concept-key supervision provides orthogonal
benefit to general-purpose pretraining but overlaps with the visual
features learned through domain-specific masked image modeling.
This suggests that SHARP is most valuable when the encoder starts from a
domain-naive initialization---a common scenario in resource-constrained
settings or when adopting new imaging modalities where domain-pretrained
models are not yet available.

\noindent\textbf{Limitations.}
SHARP requires bounding-box scene-graph annotations (weakly supervised
grounding is a natural next step), text-overlap metrics lag ImageNet-FT
on IU~X-Ray due to MIMIC-CXR phrasing bias, and the RAD-DINO
differential-LR gain does not reach significance ($p{=}0.08$).
HERGen~\cite{hergen} ($F1{=}0.317$) and MAIRA-2~\cite{maira2}
($F1{=}0.427$) exceed SHARP-FT but rely on longitudinal data or 7B
LLMs, precluding direct comparison.
We presented \textbf{SHARP}: scene-graph contrastive pretraining along
region, entity, and polarity axes achieves CheXbert F1\,=\,0.3032
($+15.8\,\%$, $p{<}0.001$); differential LR unlocks gains on
domain-pretrained RAD-DINO ($+1.7\,\%$).
Code and checkpoints will be released upon acceptance.

% ------------------------------------------------------------------
% References
% ------------------------------------------------------------------

\bibliographystyle{splncs04}

\begin{thebibliography}{99}

\bibitem{mimiccxr}
Johnson, A.E.W., et al.: MIMIC-CXR, a de-identified publicly available
database of chest radiographs with free-text reports.
Scientific Data 6(1), 317 (2019)

\bibitem{mimicextcxrqba}
M\"{u}ller, P., Jungmann, F., Kaissis, G., Rueckert, D.:
A Structured, Tagged, and Localized Visual Question Answering Dataset
with Full Sentence Answers and Scene Graphs for Chest X-ray Images.
In: ICLR (2026). \url{https://openreview.net/forum?id=LrmyW9JLYq}

\bibitem{cxrmate}
Nicolson, A., Dowling, J., Anderson, D., Koopman, B.:
Longitudinal Data and a Semantic Similarity Reward for Chest X-Ray
Report Generation.
Informatics in Medicine Unlocked 50, 101585 (2024).
\url{https://doi.org/10.1016/j.imu.2024.101585}

\bibitem{clip}
Radford, A., et al.: Learning Transferable Visual Models from Natural
Language Supervision.
In: Proceedings of the 38th ICML, pp.~8748--8763 (2021)

\bibitem{convirt}
Zhang, Y., et al.: Contrastive Learning of Medical Visual Representations
from Paired Images and Text.
In: Proceedings of Machine Learning for Health, pp.~2--25 (2022)

\bibitem{biovil}
Bannur, S., et al.: Learning to Exploit Temporal Structure for
Biomedical Vision-Language Processing.
In: Proceedings of CVPR, pp.~15016--15027 (2023)

\bibitem{medclip}
Wang, Z., et al.: MedCLIP: Contrastive Learning from Unpaired Medical
Images and Text.
In: Proceedings of EMNLP, pp.~3876--3887 (2022)

\bibitem{r2gen}
Chen, Z., et al.: Generating Radiology Reports via Memory-Driven
Transformer.
In: Proceedings of EMNLP, pp.~1439--1449 (2020)

\bibitem{rgrg}
Tanida, T., et al.: Interactive and Explainable Region-Guided Radiology
Report Generation.
In: Proceedings of CVPR, pp.~7433--7442 (2023)

\bibitem{supcon}
Khosla, P., et al.: Supervised Contrastive Learning.
In: Advances in NeurIPS 33, pp.~18661--18673 (2020)

\bibitem{vit}
Dosovitskiy, A., et al.: An Image is Worth 16x16 Words: Transformers
for Image Recognition at Scale.
In: Proceedings of ICLR (2021)

\bibitem{chexbert}
Smit, A., et al.: CheXbert: Combining Automatic Labelers and Expert
Annotations for Accurate Radiology Report Labeling Using BERT.
In: Proceedings of EMNLP, pp.~1500--1519 (2020)

\bibitem{radgraph}
Jain, S., et al.: RadGraph: Extracting Clinical Entities and Relations
from Radiology Reports.
In: Advances in NeurIPS Datasets and Benchmarks Track (2021)

\bibitem{gloria}
Huang, S.C., et al.: GLoRIA: A Multimodal Global-Local Representation
Learning Framework for Label-Efficient Medical Image Recognition.
In: Proceedings of ICCV, pp.~3942--3951 (2021)

\bibitem{hardmining}
Faghri, F., et al.: VSE++: Improving Visual-Semantic Embeddings with
Hard Negatives.
In: Proceedings of BMVC (2018)

\bibitem{harder_negatives}
Yuksekgonul, M., et al.: When and Why Vision-Language Models Behave
Like Bags-of-Words, and What to Do About It?
In: Proceedings of ICLR (2023)

\bibitem{curriculumlearning}
Bengio, Y., et al.: Curriculum Learning.
In: Proceedings of the 26th ICML, pp.~41--48 (2009)

\bibitem{bleu}
Papineni, K., Roukos, S., Ward, T., Zhu, W.J.: BLEU: A Method for
Automatic Evaluation of Machine Translation.
In: Proceedings of ACL, pp.~311--318 (2002)

\bibitem{rouge}
Lin, C.Y.: ROUGE: A Package for Automatic Evaluation of Summaries.
In: Text Summarization Branches Out, pp.~74--81 (2004)

\bibitem{meteor}
Banerjee, S., Lavie, A.: METEOR: An Automatic Metric for MT Evaluation
with Improved Correlation with Human Judgments.
In: Proceedings of the ACL Workshop on Intrinsic and Extrinsic
Evaluation Measures for MT and/or Summarization, pp.~65--72 (2005)

\bibitem{bertscore}
Zhang, T., Kishore, V., Wu, F., Weinberger, K.Q., Artzi, Y.:
BERTScore: Evaluating Text Generation with BERT.
In: Proceedings of ICLR (2020)

\bibitem{cxrbert}
Smit, A., Bhatt, D., Yu, J., Rajpurkar, P.:
CXR-BERT: A Pre-trained Language Model for Clinical Radiology Report
Summarization.
arXiv:2110.11309 (2021)

\bibitem{cmn}
Chen, Z., Shen, Y., Song, Y., Wan, X.:
Cross-modal Memory Networks for Radiology Report Generation.
In: Proceedings of ACL, pp.~5904--5914 (2021)

\bibitem{cvt2gpt2}
Nicolson, A., Dowling, J., Koopman, B.:
Improving Chest X-Ray Report Generation by Leveraging Warm Starting.
Artificial Intelligence in Medicine 144, 102633 (2023)

\bibitem{raddino}
Perez-Garcia, F., et al.: RAD-DINO: Exploring Scalable Medical Image
Encoders Beyond Text Supervision.
arXiv:2401.10815 (2024)

\bibitem{hergen}
Dalla Serra, F., et al.: HERGen: Elevating Radiology Report Generation
with Longitudinal Data.
In: Proceedings of ECCV (2024)

\bibitem{maira2}
Bannur, S., et al.: MAIRA-2: Grounded Radiology Report Generation.
arXiv:2406.04449 (2024)

\bibitem{camanet}
Wang, Z., et al.: CAMANet: Class Activation Map Guided Attention Network
for Radiology Report Generation.
IEEE Trans.\ Med.\ Imaging 42(11), 3282--3295 (2023)

\bibitem{biovilt}
Bannur, S., et al.: Learning to Exploit Temporal Structure for
Biomedical Vision-Language Processing.
In: Proceedings of CVPR, pp.~15016--15027 (2023)

\bibitem{iuxray}
Demner-Fushman, D., et al.: Preparing a collection of radiology examinations
for distribution and retrieval.
Journal of the American Medical Informatics Association 23(2), 304--310 (2016)

\bibitem{pubmedbert}
Gu, Y., Tinn, R., Cheng, H., Lucas, M., Usuyama, N., Liu, X., Naumann, T.,
Gao, J., Poon, H.: Domain-Specific Language Model Pretraining for Biomedical
Natural Language Processing.
ACM Transactions on Computing for Healthcare 3(1), 1--23 (2021)

\end{thebibliography}

\end{document}
